\section{\textbf{Solutions Through Approximations}}

To compute solutions numerically, we employ Picard Iteration using the recursive integral formula

\[ \phi_{0}(t) = x_0 \]

\[ \phi_{k+1}(t) = x_0 + \int_{t_0}^t f(\phi_k(s)) ds\]

\noindent The following figures compare the converged Picard solutions against a standard Runge-Kutta (RK45) numerical solver to verify the accuracy of the approximation.

\begin{figure}[htbp]
    \centering
    
    % --- First Row ---
    \begin{subfigure}[b]{0.24\textwidth}
        \centering
        \includegraphics[width=\textwidth]{final/figures/cp2/Picard1.png}
        \caption{Spatial Trajectory}
        \label{fig:img1}
    \end{subfigure}
    \hfill % Adds horizontal space to push the images apart
    \begin{subfigure}[b]{0.24\textwidth}
        \centering
        \includegraphics[width=\textwidth]{final/figures/cp2/Picard 2.png}
        \caption{X Components}
        \label{fig:img2}
    \end{subfigure}
    
    \label{fig:four_images}
\end{figure}

\begin{figure}[htbp]
    \centering
    % --- Second Row ---
    \begin{subfigure}[b]{0.24\textwidth}
        \centering
        \includegraphics[width=\textwidth]{final/figures/cp2/Picard3.png}
        \caption{Y Components}
        \label{fig:img3}
    \end{subfigure}
    \hfill
    \begin{subfigure}[b]{0.24\textwidth}
        \centering
        \includegraphics[width=\textwidth]{final/figures/cp2/Picard4.png}
        \caption{Orientation}
        \label{fig:img4}
    \end{subfigure}
    
    \label{fig:four_images}
\end{figure}
